We organize the experiments around three questions:
\textbf{(Q1) Heterogeneity:} Under a binary task reward, how does joint
training compare with training each task separately?
\textbf{(Q2) Scaling:} How do parallel experience and reset-state sampling
affect learning from a binary task reward?
\textbf{(Q3) Reward formulation:} How does a binary task reward compare
with dense assembly-progress shaping as heterogeneous training scales?
We additionally explore multi-pair assembly and inspect the learned policy
for behaviors that arise without direct supervision.



For heterogeneity, we compare one policy trained across all 20 ATB1 tasks
with one separately trained policy per task, using matched per-run
interaction budgets and otherwise identical settings. We use 64K environments
per run to keep the comparison affordable; the shared policy reaches
98.2\% training success at this scale.

For parallel scaling, GPUs, environments, and rollout batch size increase
together while the task distribution, reward, policy, and environment remain
fixed. The sampling ablation compares against uniform sampling over the same reset
bank instead of beta sampling.
For the scaling evaluation, all checkpoints use stochastic actions and a common
frozen mixed-reset bank, sampled uniformly within each task. Success requires
full assembly within 15 seconds without a simultaneous safety failure.
The reward ablation adds two dense assembly-progress terms to
Eq.~\eqref{eq:full-reward}: 
\begin{equation}
    r_t^{\mathrm{dense}} = r_t +
    c_2 r_t^{\mathrm{reach}} + c_3 r_t^{\mathrm{align}},
    \qquad c_2,c_3 \geq 0,
    \label{eq:reward-ablation}
\end{equation}
where $r_t^{\mathrm{reach}}$ rewards end-effector proximity to the part 
centroid and $r_t^{\mathrm{align}}$ rewards alignment with the target assembly
orientation. These are common shaping terms in manipulation and assembly RL
\cite{narang2022factory,tang2024automate,yin2026omnireset}.
We set $c_2=c_3=0.01$, matching the coefficients used in
OmniReset~\cite{yin2026omnireset}.
Across all experiments, success rates cover all reset configurations.

\begin{table*}[!t]
\centering
\renewcommand{\arraystretch}{1.2}
\caption{\textbf{Success rate for multi-pair assembly} 
Each policy is trained and evaluated on the same number of active assembly pairs, from 1--10. 
AlphaNIST solves these tasks without a prescribed order or explicit task decomposition. 
For example, when there are 10 active pairs, the policy receives a success reward only when all 10 are in their assigned board slots. 
}
\label{tab:seq-success}
\begin{tabularx}{\textwidth}{l*{10}{>{\centering\arraybackslash}X}}
\toprule
\# Active\ Pairs & 1 & 2 & 3 & 4 & 5 & 6 & 7 & 8 & 9 & 10 \\
\midrule
Success Rate (\%) & 99 & 97 & 95 & 93 & 92 & 90 & 89 & 87 & 86 & 84 \\
\bottomrule
\end{tabularx}
\end{table*}
 


\subsection{Assembly-Pair Heterogeneity Enables Learning}
From randomly initialized states, separately-trained policies have $\leq$0.3\% success on M4 nut threading, as well as 4\,mm peg insertion and rod
insertion. The main obstacle is that the small parts are
difficult to grasp reliably from the table. 
In contrast, the shared
policy produces successful grasp-and-assembly rollouts on these same tasks
without task-specific changes. This result suggests that shared experience
can help with assemblies that are difficult to learn separately.
This benefit depends strongly on the reward formulation.
At their final checkpoints, the 20 separately trained dense-reward policies
reach a mean success rate of 78.3\%, ranging from 45.0\% on
M4 threading to 93.6\% on the medium gear. Joint dense training reaches 17.8\%
overall success at 25K iterations, with low completion across tasks
(Fig.~\ref{fig:heterogeneity-ablation}).
One possible explanation
is that dense shaping terms favor different partial solutions across geometries,
whereas the binary task reward gives a unified full-assembly objective for every task.


Joint training reaches higher success at lower total cost.
At 64K parallel environments across 25K iterations, the \textbf{shared policy trained with the sparse binary task reward reaches 98.2\% training success} in 410 GPU-hours
(Fig.~\ref{fig:heterogeneity-dynamics}), while the 20 separately trained policies average 81.4\% and consume 16.1K total GPU-hours,
39.3$\times$ more device-hours.

\subsection{Sparse Binary Task Rewards Benefit from Scale}
Performance with the sparse binary task reward improves rapidly with increased parallel experience (Fig.~\ref{fig:parallel-scaling}).
With beta sampling, evaluation success rises from
38.6\% at 1K environments to 78.9\% at 4K and 86.9\% at 8K.
With dense reward shaping, success remains at 11.8--12.0\% over 1K--4K environments.
Beta sampling provides an additional gain over uniform reset sampling. At 256K environments, success with the sparse binary task reward reaches \textbf{97.4\% with beta sampling} versus 84.3\% with uniform sampling; the corresponding results with dense reward shaping are 77.4\% and 34.3\%. While the advantage of the sparse task objective is already visible at modest scale, beta sampling provides further gains as
parallel experience increases.



\subsection{Emergent Strategies over Short Horizons}



The policy exhibits a broad range of dexterous strategies across
assemblies (Fig.~\ref{fig:emergent-assembly-strategies}), including grasp adaptation, extrinsic reorientation, and contact-driven assembly. It releases and reacquires parts when the initial grasp is unsuitable for insertion, uses surrounding geometry to change object orientation, and discovers intermittent tapping and pivoting motions.
These behaviors emerge within a single policy, without demonstrations or rewards specifying the individual
maneuvers.
For a peg, a successful pickup may not produce a grasp suitable for insertion; the
policy drops a center-held peg and regrasps it near one end, exposing enough
shaft length for complete insertion. For a nut, repeated taps on one side engage
and advance the thread, using intermittent contact to make assembly progress.

\subsection{Emergent Strategies over Long Horizons}
Beyond the main single-assembly experiments, we train policies with multiple
assembly pairs per episode. The binary task reward is issued only when every
active pair is assembled; episodes allow 15 seconds per pair and terminate upon success. 
Table \ref{tab:seq-success} shows the success rates for episodes with up to ten active pairs.
This multi-pair setting reveals coordination over substantially longer horizons (Fig.~\ref{fig:long-horizon-multi-asset}). The policy need not finish
assembling a pair before turning to another, and its approach to
one assembly changes based on neighboring parts. 
These planning-like behaviors---taking actions whose utility is realized later and adapting one assembly to the state of another---emerge without an explicit planner, policy switching, or prescribed sequence.


